\documentclass[a4paper,openany]{book}

% 文章排版
\usepackage{indentfirst}
\setlength{\parindent}{0em}
\usepackage{autobreak}
\allowdisplaybreaks
\usepackage{parskip}
\usepackage{geometry}
\geometry{top=3cm,bottom=3cm,left=2cm,right=2cm}
\usepackage{anyfontsize}
\usepackage{hyperref}

% 数学物理宏包
\usepackage{amsmath}
\usepackage{esint}
\usepackage{tabularray}
\UseTblrLibrary{amsmath}
\usepackage{physics}
\renewcommand\div\divisionsymbol

% 设置数学字体

% Linux Libertine 风格设置
% 不需要 fontspec 和 unicode-math
\usepackage[libertine,vvarbb]{newtx}
\usepackage{bm}
\renewcommand\mathbf\bm

% 中文
\usepackage[fontset=none]{ctex}
\setCJKmainfont{SimSun}[BoldFont = Source Han Sans CN Medium, ItalicFont=KaiTi]

% 符号重定义
\AtBeginDocument{
    \let\uglyepsilon\epsilon
    \renewcommand\epsilon\varepsilon
    \let\uglyOmega\Omega
    \renewcommand\Omega\varOmega
    \let\uglyleq\leq
    \renewcommand\leq\leqslant
    \let\uglygeq\geq
    \renewcommand\geq\geqslant
}

\begin{document}

\title{\textbf{宇宙学讲义}}
\author{北京大学\ \ \ \ 张植竣}
\date{\today}
\maketitle

\tableofcontents

\chapter{宇宙学原理}

\section{宇宙学原理与Robertson-Walker度规}

宇宙学中最重要的一个原理是\textbf{宇宙学原理}：\textbf{宇宙在大尺度上是均匀和各向同性的}。均匀（homogeneous）指的是宇宙中各点都是平等的，没有何处是特殊的；各向同性（isotropic）指的是相对于某一点，宇宙在各个方向上看起来都是一样的。均匀的不一定是各向同性的，各向同性的也不一定是均匀的。但是如果宇宙对任意一点都是各向同性的，那么它一定是均匀的。对于人类，我们只能从所在的地球乃至太阳系观测宇宙，可以得出宇宙对于我们是各向同性的结论，但如果我们接受\textbf{哥白尼原理}：\textbf{人类在观测宇宙时并非处于一个特殊地位}。即在宇宙的任何一点观测宇宙都和地球上观测的结果相同，则宇宙均匀和各向同性都随之成立。

% Figure 1: homogeneous vs isotropic

宇宙在多大尺度上符合宇宙学原理呢？典型星系团的尺度是$10^1\,\mathrm{Mpc}$的量级，超星系团可以达到$10^2\,\mathrm{Mpc}$，可观测宇宙的直径大约是$4.4\,\mathrm{Gpc}$。在数百个$\mathrm{Mpc}$的尺度上，宇宙看起来是光滑的，没有来自于局部引力塌陷形成的结构。例如，如果我们在天球相对的两端各取一个边长为$100\,\mathrm{Mpc}$的立方体，并计算其中物质的密度，两者应该给出极其相近的结果。宇宙微波背景辐射（CMB）的涨落甚至只有$10^{-6}\,\mathrm{K}$的量级。

在宇宙学中，需要定义一种特别的理想观测者：基本观测者。这些观测者相对整个宇宙流体没有相对运动。宇宙流体由宇宙中所有物质（包括暗物质）平均后的运动来定义。简单地说，他们不受万有引力的轨道运动影响，只随着整个宇宙一起膨胀。

我们首先要引入“度规”的概念。对于相隔无穷小的两个点，可以认为它们是直线相连，定义线元
\begin{equation}
    \dd{s}^2 = \sum_{i,j}g_{ij} \dd{x}^i \dd{x}^j,\quad x^i = t,\ x,\ y,\ z
\end{equation}
对于狭义相对论下的Minkowski时空：
\begin{equation}
    \dd{s}^2 = -(c\dd{t})^2 + \dd{x}^2 + \dd{y}^2 + \dd{z}^2
\end{equation}
根据宇宙学原理，宇宙空间的度规只能有如下的形式，称为\textbf{Robertson-Walker度规}，简称RW度规：
\begin{equation}
    \dd{s}^2 = -(c\dd{t})^2 + a^2(t) \left[\frac{\dd{r}^2}{1 - kr^2} + r^2\left(\dd{\theta}^2 + \sin^2\theta \dd{\phi}^2\right)\right]
\end{equation}
其中，$a(t)$称为尺度因子，$k$为归一化的宇宙曲率。

尺度因子$a(t)$描述了宇宙的膨胀率，我们可以定义：
\begin{equation}
    \mathbf{r} = a(t)\mathbf{x}
\end{equation}
其中$\mathbf{r}$称为物理坐标，而$\mathbf{x}$称为共动坐标。记现在的宇宙时间是$t_0$，规定当前的共动坐标等于测量到的物理坐标，即尺度因子$a(t_0) \equiv 1$。一般认为宇宙始终不断膨胀，则过去时刻的尺度因子$a(t)$总是小于1。我们可以这样来理解这个关系：想象宇宙的空间部分是在一个气球上面，星系可以看作球面上的点，而宇宙的演化类似于气球的膨胀，尺度因子代表气球的半径。对球面而言，球面上的点的相对位置是不变的，这就是共动坐标的含义。但对于外部观测者而言，整个气球是在膨胀的，气球上点的相对位置变大了。

% Figure 2, the balloon model of expanding universe

宇宙曲率$k$描述了宇宙的形状。$k$的取值有以下三种情况：
\begin{enumerate}
    \item $k = 0$：宇宙为一个平坦空间，最终会膨胀到一个最大半径，但膨胀速度趋于零；
    \item $k = 1$：宇宙为一个球形空间，闭的，最终会塌缩到一点；
    \item $k = -1$：宇宙为一个双曲空间，开的，最终会无限加速膨胀。
\end{enumerate}

\section{宇宙学红移与Hubble定律}

红移的定义是：
\begin{equation}
    z \equiv \frac{\lambda}{\lambda_0} - 1
\end{equation}
其中$\lambda$为观测者接收到的波长，$\lambda_0$为光谱的实验室静止波长。在实际观测中，光子发射时的波长一般为静止波长，则有：
\begin{equation}
    z = \frac{\lambda_o}{\lambda_e} - 1
\end{equation}
下标$o$表示观测（observasion），$e$表示发射（emission）。随着宇宙膨胀，光子的波长也会随之膨胀：
\begin{equation}
    z = \frac{\lambda_o}{\lambda_e} - 1 = \frac{a(t_o)}{a(t_e)} - 1
\end{equation}
广义相对论对光子的严谨计算表明，这个结果不依赖于宇宙的膨胀历史，无论尺度因子的变化快慢如何，最终的结果只依赖于尺度因子在发射时刻和接收时刻的值。

如果在当前时刻$t_0$我们接收到从遥远星系发出的光子，而发射光子时的宇宙时间为$t$。对于当前的观测者，尺度因子$a(t_0) = 1$，则有：
\begin{equation}
    z = \frac{1}{a(t)} - 1
\end{equation}

对于在过去不久的时刻发射的光子，$t = t_0 - \Delta t$，其中$\Delta t \ll t_0$，可以对上式做Taylor展开：
\begin{equation}
    1 + z = \frac{1}{a(t_0 - \Delta t)} = a(t_0) + \Delta t \dot{a}(t_0)
\end{equation}
我们引入一个参数，称为\textbf{Hubble常数}：
\begin{equation}
    H(t) \equiv \frac{\dot{a}(t)}{a(t)}
\end{equation}
Hubble常数实际上并非“常数”，它的值也会随着宇宙演化而改变，其当前值记为：
\begin{equation}
    H_0 = H(t_0) = \dot{a}(t_0)
\end{equation}
将(1.10)式代入(1.9)式，并注意到$a(t_0) = 1$，可得：
\begin{equation}
    1 + z = 1 + H_0 \Delta t
\end{equation}
即:
\begin{equation}
    z = H_0 \Delta t
\end{equation}
如果光子需要$\Delta t$时间到达观测者，则光子发射处距离观测者为$d = c\Delta t$，又因为径向的退行速度（按经典Doppler红移的定义）$v_r = cz$，可以得到\textbf{Hubble定律}：
\begin{equation}
    v_r = H_0 \Delta d
\end{equation}
注意，这仅在红移较小（$z \ll 1$）时成立。

根据Planck卫星在2018年发表的测量结果：
\begin{equation}
    H_0 = 67.66\,\mathrm{km \cdot s^{-1} \cdot Mpc^{-1}}
\end{equation}
可以发现，Hubble常数实际上具有一个时间倒数的量纲（两个距离的量纲相消），于是有：
\begin{equation}
    H_0 = 2.193 \times 10^{-18}\,\mathrm{s^{-1}}
\end{equation}
由于测量的不确定性，实际计算中会将它参数化为：
\begin{equation}
    H_0 = 100h \,\mathrm{km \cdot s^{-1} \cdot Mpc^{-1}}
\end{equation}
则无量纲参数$h$如今的取值是$h = 0.6766$。

通过哈勃常数可以给出宇宙年龄和尺度的典型值。宇宙年龄近似为\textbf{Hubble时间}：
\begin{equation}
    t_H = H_0^{-1} = 4.56 \times 10^{17}\,\mathrm{s} = 14.45 \,\mathrm{Gyr}
\end{equation}
可观测宇宙的尺度近似为\textbf{Hubble长度}：
\begin{equation}
    d_H = ct_H = 1.37 \times 10^{26}\,\mathrm{m} = 4.43\,\mathrm{Gpc}
\end{equation}
实际宇宙的年龄和长度与宇宙学模型（即宇宙各组分的比例）有关，在第三章会详细讲解。它们与Hubble时间和Hubble长度大多只相差一个约等于1的常数。

\chapter{宇宙的演化}

\section{Friedmann方程与流体方程}

宇宙的演化由\textbf{Einstein方程}确定：
\begin{equation}
    G_{\mu\nu} = \frac{8\pi G}{c^4}T_{\mu\nu}
\end{equation}
方程的左边$G_{\mu\nu}$是Einstein张量，表征宇宙空间的弯曲；方程的右边$T_{\mu\nu}$是能动张量，表征空间中物质与能量的分布；$G$为万有引力常数，透过弱场近似以及慢速近似，可以从爱因斯坦场方程退化为牛顿重力定律。

实际上我们可以在方程的左边加上一个与度规张量$g_{\mu\nu}$成正比的项，这样就得到方程：
\begin{equation}
    G_{\mu\nu} + \Lambda g_{\mu\nu} = \frac{8\pi G}{c^4}T_{\mu\nu}
\end{equation}
这个比例系数$\Lambda$就称为\textbf{宇宙学常数}，现在一般认为这个常数表示宇宙中的暗能量，或称为真空能。

宇宙尺度因子的演化取决于物质组分，不同的物质组分具有不同的物态方程：
\begin{equation}
    p = w\rho
\end{equation}
其中$p$为压强，$\rho$为密度，$w$是一个与时间无关的常数，称为物态参数。物质组分可以分为三类：
\begin{enumerate}
    \item $p_m = 0$，非相对论性的物质，一般简称为物质（或尘埃），包含重子物质和暗物质；
    \item $p_r = \dfrac{1}{3}\rho_r$，相对论性的物质，即辐射，包含光子、中微子等；
    \item $p_\Lambda = -\rho_\Lambda$，宇宙学常数。
\end{enumerate}
在RW度规下，将物态方程代入Einstein方程(1.17)式可以得到：
\begin{align}
    \left(\frac{\dot{a}}{a}\right)^2 &= \frac{8\pi G}{3}\rho - \frac{kc^2}{a^2} + \frac{\Lambda c^2}{3} \\
    2 \frac{\ddot{a}}{a} + \left(\frac{\dot{a}}{a}\right)^2 &= - \frac{kc^2}{a^2} -\frac{8\pi G}{c^2}\rho + \Lambda c^2
\end{align}
这两个就是著名的\textbf{Friedmann方程}，由俄国物理学家Alexander Friedmann发表于1922年。它们互相独立，共同描述了一个均匀且各向同性的膨胀宇宙模型。

另一个重要的方程是所谓\textbf{流体方程}：
\begin{equation}
    \dot{\rho} + 3\frac{\dot{a}}{a}\left(\rho + p\right) = 0
\end{equation}
这个方程可以从热力学第一定律$\mathrm{d}E = -P\mathrm{d}V + T\mathrm{d}S$中令$\mathrm{d}S = 0$得到，即宇宙是绝热膨胀的，和“宇宙外部”没有能量交换。

\section{物质组分}

首先，我们简要地讨论一下宇宙中的物质组分。通过流体方程(2.6)式，我们可以推导出不同的物质组分随尺度因子的演化。

\begin{enumerate}
    \item 对于物质，$p_m = 0$，代入(2.6)式得：
    \begin{equation}
        \dot{\rho}_m + 3\frac{\dot{a}}{a}\left(\rho_m + 0\right) = 0
    \end{equation}
    即：
    \begin{equation}
        \frac{\dot{\rho}_m}{\rho_m} = -3\frac{\dot{a}}{a}
    \end{equation}
    积分得：
    \begin{equation}
        \rho_m \propto a^{-3}
    \end{equation}
    即：
    \begin{equation}
        \rho_m(z) = \rho_{m,0}(1 + z)^3
    \end{equation}

    \item 对于辐射，$p_r = \dfrac{1}{3}\rho_r$，代入(2.6)式得：
    \begin{equation}
        \dot{\rho}_r + 3\frac{\dot{a}}{a}\left(\rho_r + \frac{1}{3}\rho_r\right) = 0
    \end{equation}
    即：
    \begin{equation}
        \frac{\dot{\rho}_r}{\rho_r} = -4\frac{\dot{a}}{a}
    \end{equation}
    积分得：
    \begin{equation}
        \rho_r \propto a^{-4}
    \end{equation}
    即：
    \begin{equation}
        \rho_r(z) = \rho_{r,0}(1 + z)^4
    \end{equation}
    由于辐射的衰减最快，在如今的宇宙中光子辐射的比重只有约$5.4 \times 10^{-5}$，中微子则占约$1.44 \times 10^{-3}$。一般在宇宙学演化的计算中可以忽略辐射的影响。

    值得一提的是，宇宙中辐射能量密度的主要贡献（超过$99\%$）是宇宙微波背景辐射（CMB）。CMB光子能量分布在宇宙演化中保持黑体辐射的形式，根据Stefan-Boltzmann定律：
    \begin{equation}
        \rho_r \propto T^4
    \end{equation}
    则宇宙温度（即CMB的温度）和红移的关系为：
    \begin{equation}
        T(z) = T_0(1 + z)
    \end{equation}

    \item 对于宇宙学常数，$p_\Lambda = -\rho_\Lambda$，代入(2.6)式得：
    \begin{equation}
        \dot{\rho}_\Lambda + 3\frac{\dot{a}}{a}\left(\rho_\Lambda - \rho_\Lambda\right) = 0
    \end{equation}
    即：
    \begin{equation}
        \dot{\rho}_\Lambda = 0
    \end{equation}
    积分得：
    \begin{equation}
        \rho_\Lambda = \mathrm{const.}
    \end{equation}
    即宇宙学常数的能量密度并不随尺度因子而变化。
\end{enumerate}

在Friedmann方程(2.4)中，注意到方程左边实际上就是$H^2$，可以将方程化为：
\begin{equation}
    \frac{8\pi G}{3H^2}\rho - \frac{kc^2}{a^2 H^2} + \frac{\Lambda c^2}{3H^2} = 1
\end{equation}
定义\textbf{临界密度}：
\begin{equation}
    \rho_c = \frac{3H^2}{8\pi G}
\end{equation}
可以定义所谓的密度参数：
\begin{align}
    \Omega_m &= \frac{\rho}{\rho_c} = \frac{8\pi G}{3H^2}\rho,\quad \rho = \rho_m + \rho_r \approx \rho_m \\
    \Omega_k &= \frac{\rho_k}{\rho_c} = -\frac{kc^2}{a^2 H^2},\quad \rho_k = -\frac{3kc^2}{8\pi Ga^2} \\
    \Omega_\Lambda &= \frac{\rho_\Lambda}{\rho_c} = \frac{\Lambda c^2}{3H^2},\quad \rho_\Lambda = -\frac{\Lambda c^2}{8\pi G} \\
\end{align}
则(2.7)式化为：
\begin{equation}
    \Omega_m + \Omega_k + \Omega_\Lambda = 1
\end{equation}
根据Planck卫星在2018年发表的测量结果，这三个密度参数的值为：
\begin{equation}
    \Omega_m = 0.30966,\quad \Omega_k = 0,\quad \Omega_\Lambda = 0.68884
\end{equation}
重子物质、暗物质、光子和中微子的密度参数为：
\begin{equation}
    \Omega_b = 0.04897,\quad \Omega_{dm} = 0.26069,\quad \Omega_\gamma = 5.402 \times 10^{-5},\quad \Omega_\nu = 0.00144
\end{equation}
临界密度的值为：
\begin{equation}
    \rho_c = 8.5988 \times 10^{-27}\,\mathrm{kg/m^3}
\end{equation}

\section{Hubble常数的演化}

从Friedmann方程可以看出，Hubble常数会随着时间演化，演化规律为：
\begin{equation}
    \left(\frac{H(z)}{H_0}\right)^2 = \Omega_{m,0}(1 + z)^3 + \Omega_{k,0} (1 + z)^2 + \Omega_{\Lambda,0}
\end{equation}

\chapter{宇宙学模型}

最简单的情况是$k = \Lambda = 0$，即物质主导的、没有暗能量的平坦宇宙。这被称为\textbf{Einstein-de Sitter模型}，此时Friedmann方程(2.4)变为：
\begin{equation}
    \left(\frac{\dot{a}}{a}\right)^2 = \frac{8\pi G}{3}\rho
\end{equation}
注意物质密度遵循$\rho = \rho_0 a^{-3}$，则有：
\begin{equation}
    \left(\frac{\dot{a}}{a}\right)^2 = \frac{8\pi G}{3}\rho_0 a^{-3}
\end{equation}
由于$\rho_k = \rho_\Lambda = 0$当前的物质密度$\rho_0$还需满足：
\begin{equation}
    \rho_0 = \rho_c = \frac{3H_0^2}{8\pi G}
\end{equation}
则(3.3)式化简为：
\begin{equation}
    \dot{a} = \dv{a}{t} = \frac{H_0}{\sqrt{a}}
\end{equation}
最后积分得：
\begin{equation}
    a = \left(\frac{3}{2}H_0 t\right)^{\tfrac{2}{3}}
\end{equation}
这便是尺度因子随时间的变化关系。

对于一个辐射主导的宇宙，同理有：
\begin{equation}
    \left(\frac{\dot{a}}{a}\right)^2 = \frac{8\pi G}{3}\rho_{r,0} a^{-4}
\end{equation}
则：
\begin{equation}
    a^2 H = \sqrt{\frac{8\pi G}{3}\rho_{r,0}} = \sqrt{H_0^2 \Omega_{r,0}}
\end{equation}

当前对宇宙的观测支持一个$k = 0$且$\Lambda > 0$的宇宙，可以计算得出它满足：
\begin{equation}
    a^3 = \frac{\Omega_{m,0}}{2\Omega_{\Lambda,0}}[\cosh(\sqrt{3\Lambda}t) - 1]
\end{equation}

\chapter{距离与红移}

在宇宙中距离的定义和测量容易让人迷惑。在RW度规中，$r$给出的是共动距离，而物理距离$d = a(t)r$。然而物理距离具有操作性上的困难，无法直接测量，因此需要引进一些更具有操作性的距离概念。

\section{光度距离}

在静止的Euclid时空中，考虑一个距离我们$d$的光源，具有光度$L$，我们收到的光通量为：
\begin{equation}
    F = \frac{L}{4\pi d^2}
\end{equation}
因此如果知道光源的光度，利用测到的光通量，我们可以得到距离：
\begin{equation}
    d = \sqrt{\frac{L}{4\pi F}}
\end{equation}
在RW度规下，考虑一个在固定$d$处的光源，其光度为$L(t)$，它在宇宙时刻$t$处发光，并在$t_0$时刻被我们观测到。假设光源在时间段$\Delta t_e$内发出的光，波长为$\lambda_e$，在时间段$\Delta t_o$内被接收到，此时波长为$\lambda_o$。由于：
\begin{equation}
    \frac{\lambda_e}{\lambda_o} = \frac{\Delta t_e}{\Delta t_o} = \frac{a(t)}{a(t_0)} = \frac{1}{1 + z}
\end{equation}
光度和亮度本质衡量的是光的功率，即单位时间发射或接收的能量，而这两者分别为：
\begin{equation}
    P_e = \frac{hc}{\lambda_e\Delta t_e},\quad P_o = \frac{hc}{\lambda_o\Delta t_o}
\end{equation}
两者之比为：
\begin{equation}
    \frac{P_o}{P_e} = \left(\frac{a(t)}{a(t_0)}\right)^2 = \frac{1}{(1 + z)^2}
\end{equation}
即：
\begin{equation}
    F = \frac{L}{4\pi d^2} \cdot \frac{1}{(1 + z)^2}
\end{equation}
定义\textbf{光度距离}：
\begin{equation}
    d_L = d(1 + z)
\end{equation}
则：
\begin{equation}
    F = \frac{L}{4\pi d_L^2}
\end{equation}

\section{角直径距离}

假设在距离$d$处有一个较大的光源（如星系团）直径为$D$，它在宇宙时刻$t$处发光，并在$t_0$时刻被我们观测到。观测到它在天上的角直径为$\theta$，则按照RW度规，光源直径应表示为：
\begin{equation}
    D = a(t)d\theta
\end{equation}
即距离为：
\begin{equation}
    d = \frac{D}{a(t)\theta}
\end{equation}
在静止的Euclid时空中，距离应为：
\begin{equation}
    d = \frac{D}{\theta}
\end{equation}
则可以定义\textbf{角直径距离}：
\begin{equation}
    d_A = \frac{D}{\theta} = \frac{d}{1 + z}
\end{equation}

% \section{视界}

\chapter{黑洞}

不旋转且不带电的黑洞称为\textbf{史瓦西黑洞}，其视界半径为：
\begin{equation}
    r_S = \frac{2GM}{c^2}
\end{equation}
所有进入视界的物质，包括光子，都不可能从视界出来。

忽略掉$2\pi$等等的因子，普朗克质量的意义大约是一个史瓦西半径等同于康普顿波长的黑洞所带有的质量。这给出了微观粒子的质量上限。
\begin{equation}
    \frac{2Gm}{c^2} = \frac{2\pi\hbar}{mc}
\end{equation}
忽略数学常数可以化为：
\begin{equation}
    \frac{Gm}{c^2} = \frac{\hbar}{mc}
\end{equation}
解得：
\begin{equation}
    m_P = \sqrt{\frac{\hbar c}{G}} = 21.765\,\mu\mathrm{g}
\end{equation}
这黑洞的半径大约是普朗克长度：
\begin{equation}
    l_P = \sqrt{\frac{\hbar}{Gc^3}} = 1.616 \times 10^{-35}\,\mathrm{m}
\end{equation}
普朗克时间是普朗克长度除以光速：
\begin{equation}
    t_P = \sqrt{\frac{\hbar}{Gc^5}} = 5.391 \times 10^{-44}\,\mathrm{s}
\end{equation}

\end{document}